% Summary
\emph{Sammanfattning:}
Detta kapitel introducerar objektorienterad programmering i Python genom 
konceptet klasser som sammansatt datatyp.
Vi börjar med att jämföra två tillvägagångssätt för att lösa samma problem:
ett med nuvarande förkunskaper och ett med klasser.
Denna jämförelse illustrerar fördelarna med klasser för att organisera och 
strukturera kod.

Kapitlet täcker de grundläggande koncepten klasser, objekt, attribut och 
metoder, inklusive specialmetoder (dundermetoder, 
\enquote{\foreignlanguage{english}{dunder methods}}) som tillåter anpassat 
beteende för standardoperationer i Python.
Slutligen utforskar vi arv, och visar hur man skapar nya klasser baserade på 
befintliga för att utöka funktionalitet utan kodduplikation.

% Motivation and intended learning outcomes
\emph{Avsedda lärandemål:}
Efter att ha genomfört detta kapitel ska du kunna:
\begin{enumerate}
  \item[LO1] Förklara skillnaden mellan en klass och ett objekt.
  \item[LO2] Skapa klasser med attribut och metoder.
  \item[LO3] Förstå och använda parametern \mintinline{python}{self}.
  \item[LO4] Implementera dundermetoder som
    \begin{itemize}
      \item \mintinline{python}{__init__},
      \item \mintinline{python}{__str__},
      \item \mintinline{python}{__eq__} och \mintinline{python}{__lt__}.
    \end{itemize}
  \item[LO5] Använda inkapsling för att skydda data.
\end{enumerate}
Utöver detta har du även möjlighet att lära dig om arv genom att:
\begin{enumerate}
  \item[LO6] Skapa underklasser med hjälp av arv.
  \item[LO7] Överlagra metoder i underklasser och använda
    \mintinline{python}{super()} för att anropa föräldraklassens metoder.
  \item[LO8] Avgöra när man ska använda arv kontra komposition.
\end{enumerate}

% Prerequisites
\emph{Förkunskaper:}
För att förstå detta kapitel bör du vara bekväm med:
\begin{itemize}
  \item Grundläggande Python-syntax (variabler, funktioner, kontrollstrukturer)
  \item Dictionaries och hur man arbetar med dem
  \item Funktioner och parametrar
  \item Strängmanipulering
\end{itemize}
